\hspace*{1.6cm}With Thailand entering an aged society, loneliness and social isolation among the elderly have become critical issues. Existing monitoring solutions often lack interactivity or are too complex for elderly users. This project proposes "Buddy," an interactive companion system designed to enhance the quality of life for the elderly through natural communication and emotional monitoring. The system consists of a tablet-based application for the elderly and a mobile application for caretakers, facilitating seamless interaction. The technical architecture implements a comprehensive Speech-to-Speech (STS) pipeline, utilizing Faster Whisper for high-accuracy Thai speech recognition and a local Large Language Model (GEMMA 3:4b) for generating empathetic, persona-conditioned responses. A separate GEMMA 3:1b model analyzes conversation logs to track emotional trends, while ZegoCloud SDK powers reliable real-time video calls.

\hspace*{1.6cm}Experimental results validate the system's efficacy across multiple dimensions. The speech recognition module achieved a Word Error Rate (WER) of 10.50\% and Character Error Rate (CER) of 7.20\%, outperforming the Thonburian benchmark. Emotion classification demonstrated high reliability with an average F1-score of 0.93. Although the lightweight local LLM scored lower than cloud-based counterparts in general knowledge, it proved highly effective for emotional support and daily conversation, achieving a human Mean Opinion Score (MOS) of 4.2 out of 5. The telecommunication system showed robust performance with a 98.5\% connection success rate and an average latency of 250 ms. User surveys revealed a satisfaction score of 5.5 out of 7, suggesting that "Buddy" successfully reduces technological barriers and provides valuable emotional insights to caretakers, effectively serving as a bridge between the elderly and their families.